\section{Evaluation Protocol}
\label{sec:protocol}

An edit can change too much outside the requested region or too little inside it.
Our automated protocol records both directions, but it is a systems diagnostic:
it does not establish whether an edit is anatomically or clinically plausible.

\subsection{Identity}

We compute the cosine similarity between ArcFace~\cite{deng2019arcface}
embeddings of the largest detected face in the input and in the edit. We plot
$0.6$ as a heuristic reference threshold, not a clinically calibrated acceptance
criterion; verification thresholds depend on implementation, preprocessing, and
the desired false-match rate. A high score can also mean the model did not perform
the edit, so identity is read with the regional-change scores below.

As Agarwal and Bhrany~\cite{agarwal2026envisage} note, a full-face identity score
on a composited output is partly determined by pixels copied from the input. We
therefore treat prompt-only identity as a model-level diagnostic and composite
identity as a system-level property. Neither score measures surgical accuracy.

\subsection{Regional change and localization}

After aligning the edit onto the input, we convert both to CIELAB and take the
mean $\Delta E_{76}$ color difference inside the requested target region
($\Delta E_{\mathrm{tgt}}$) and inside a disjoint facial keep zone
($\Delta E_{\mathrm{off}}$, the upper-face region minus the target). We use the
simple 1976 distance rather than the perceptually refined
CIEDE2000~\cite{sharma2005ciede2000}; it is a reproducible pixel-change
diagnostic, not a perceptual or anatomical accuracy score. The summary statistic
is
\[
  \mathrm{loc} = \frac{\Delta E_{\mathrm{tgt}}}
                      {\Delta E_{\mathrm{tgt}} + \Delta E_{\mathrm{off}}}.
\]
The ratio is read with its components: low $\Delta E_{\mathrm{tgt}}$ identifies
a near-copy, while high $\Delta E_{\mathrm{off}}$ indicates leakage into the
measured keep zone. Related preservation metrics appear in general-domain editing
benchmarks~\cite{qian2025giebench,alebench2024,sheynin2024emuedit}.

The ratio has three limitations. First, the keep zone does not cover the entire
image; lower-face, hair, clothing, or background changes may be missed,
particularly for rhinoplasty. Second, the frontal-mesh metric abstains on profile
inputs. Third, compositing copies off-mask pixels by design, so a high ratio on
that rung verifies the implementation's preservation property rather than the
editor's quality.

\subsection{Postoperative identity reference}

For faces with a postoperative photograph, we report the ArcFace cosine between
the edit and that photograph, and beside it a per-face baseline: the cosine
between the unedited input and the same postoperative photograph, computed with
the same embedding pipeline. The baseline is what an edit has to beat before any
claim of moving closer to the observed outcome can be made in embedding space.
Both scores are identity-continuity references only. The before and after
sessions differ in pose, expression, makeup, and hair, and movement in an
identity embedding is not evidence of anatomical accuracy, so we make no
outcome-accuracy claim from these scores.

\subsection{Descriptive uncertainty}

The face, not the generated edit, is the sampling unit. Cell summaries are
medians and ranges because each model--control cell contains fewer than ten faces.
For the primary prompt-only versus masked-composite contrast, we compute paired
differences within face, procedure, and model, then resample faces (10{,}000
replicates, fixed analysis seed) to obtain a percentile interval for the median.
This interval describes the observed face set; it does not repair the small sample
or support population-level clinical inference.
